% APPENDIX. Insert after \appendix, before app:numerical.
\section{Proofs for the corner and persistence results}
\label{app:stage1-proofs}

\subsection{Tightness of the production floor}
\label{app:floor-proof}
\begin{proof}[Proof of \cref{cor:production-growth}]
The supporting-line inequality for a concave function, based at $B$,
reads $\ell(v)\leq\ell(B)+\ell'(B)(v-B)$ for $0\leq v\leq B$.
At $v=0$ it gives $B\ell'(B)\leq\ell(B)$, proving $\Gamma(B)\geq B$.
If equality holds, concavity also gives
$\ell(v)\geq(v/B)\ell(B)=v\ell'(B)$. The supporting line is
$v\ell'(B)$, so both inequalities are equalities throughout $[0,B]$.
Conversely a linear function on that interval plainly gives equality.
No assertion about production beyond $B$ is needed.

At $\eta=1$, \cref{eq:hub} becomes
\[
 p_j\geq S_j\Gamma(B)+\frac{r_j+d_0S_j}{a\ell'(B)}.
\]
The second term is strictly positive whenever $S_j>0$, since $d_0>0$
and $a\ell'(B)>0$. Thus $p_j>S_jB$ is necessary in every such
direction for any admissible concave production function.
If those inequalities hold, choose $\ell(v)=v$.
All denominators $p_j-S_jB$ are positive, including $p_j$ in a
zero-slope direction. The finite choice
\[
 a>\max_j\frac{r_j+d_0S_j}{p_j-S_jB}
\]
satisfies every corner inequality strictly. A direction with
$r_j=S_j=0$ contributes zero to this maximum. If all directions are
trivial, any positive $a$ works. This proves necessity and sufficiency
of the claimed existence condition, including zero residual slopes.
Along an unbounded-budget path, a fixed positive slope thus requires
$p_j(t)>S_jB(t)$. If $p_j(t)\leq M$, then
$S_j(t)<M/B(t)$ at positive-slope times, and hence
$S_j(t)=O(B(t)^{-1})$ including zero-slope times. These are necessary
rate restrictions; the remaining productivity term is still required.
\end{proof}

\subsection{Growing budgets and the original special cases}
\label{app:growth-proof}
\begin{proof}[Proof of \cref{prop:growth-trichotomy}]
Multiplying \cref{eq:hub} by $\Gamma$ gives the required price
\[
 P_{j,\eta}(\Omega,B)=\Gamma(B)
 \left[\frac{r_j}{\Omega}
       +S_j\frac{(d_0+\Omega)^\eta}{\Omega}\right].
\]
For $\eta=1$, the second quotient is $1+d_0/\Omega>1$.
Hence a passing direction with $p_j\leq M$ and $S_j\geq s$ requires
$M>s\Gamma(B)\geq sB$, which eventually fails.
For $\eta>1$, differentiation gives
\[
 \frac{\dd}{\dd\Omega}\frac{(d_0+\Omega)^\eta}{\Omega}
 =\frac{(d_0+\Omega)^{\eta-1}}{\Omega^2}
       [(\eta-1)\Omega-d_0].
\]
The quotient diverges at both endpoints of $(0,\infty)$, decreases
until $\Omega=d_0/(\eta-1)$, and increases thereafter. Substitution
gives \cref{eq:residual-minimum}. Consequently the required price is
at least $sB$ times a fixed positive number and again eventually exceeds
$M$. The argument applies separately at each time and imposes no
regularity or growth restriction on $a(t)$.

For $0\leq\eta<1$, $(d_0+\Omega)^\eta\geq\Omega^\eta$.
A passing direction therefore requires
$M\geq s\Gamma\Omega^{\eta-1}$, which is exactly
\cref{eq:necessary-output-growth}. For sufficiency, fix the parameters
as in the proposition, let $0<\eta<1$, and suppose
$\Omega\geq c_\Omega\Gamma^{1/(1-\eta)}$. Since
$\Gamma\geq B\to\infty$, also $\Omega\to\infty$, and
\begin{align*}
 P_{j,\eta}
 &=\frac{r_j\Gamma}{\Omega}
   +S_j\Gamma\Omega^{\eta-1}(1+d_0/\Omega)^\eta\\
 &\leq\frac{r_j}{c_\Omega}\Gamma^{-\eta/(1-\eta)}
   +S_jc_\Omega^{\eta-1}(1+d_0/\Omega)^\eta.
\end{align*}
The first term tends to zero. The limsup of the second is strictly
less than $p_j$ by the stated strict bound on $c_\Omega$; when $S_j=0$
both terms tend to zero. Finitely many directions admit a common time
beyond which all inequalities are strict. For $\eta=0$ the same exact
price is $(r_j+S_j)\Gamma/\Omega$, proving
\cref{eq:linear-value-growth}. Finally, $\Gamma=B$ and $\Omega=aB$
for linear production give the asserted rate for $a$.
\end{proof}

The sufficiency assumptions cannot be replaced by an upper price bound
and a lower residual bound alone. For example, take $\eta=1/2$,
$\ell(z)=z$, $d_0=1$, $r=0$, $S=1$, $B(t)=t$, $a(t)=10t$,
and $p(t)=\exp(-t)$. This productivity has the proposed power-law rate,
but its required price tends to $1/\sqrt{10}$ while its actual price
tends to zero. Likewise equality in the asymptotic coefficient is not
generally enough: with $\ell(z)=z$, $S=p=d_0=1$, $r=0$, and
$a=B^{\eta/(1-\eta)}$, the required price is
$(1+1/\Omega)^\eta>1$ for every finite $B$ and $0<\eta<1$.

For clarity, the older reductions are algebraic, not additional
assumptions. Substituting $\eta=1$ gives
\cref{eq:general-production}; additionally setting
$\ell(B)=\ln(1+B)$ gives $H_j,A_j$ and \cref{eq:exact}.
Keeping arbitrary $\eta$ with this production function gives
\cref{eq:iso-corner}. Applying the logarithmic formula at each time
gives \cref{eq:pathwise}; its zero-$A_j$ case is exactly
$r_j=S_j=0$. The general production floor grows at least linearly,
whereas the logarithmic special case is $(1+B)\ln(1+B)\sim B\ln B$.

The original pure-computation solution also follows. If all residual
weights vanish and external computation is affine, any strictly
increasing $F$ has the same optimizer as $C$. At a fixed local
allocation $z$, remaining resources are assigned to a channel with
$R_*=\max(r_n/p_n,r_h/p_h)$. For $R_*>0$, the reduced derivative is
$a/(1+z)-R_*$ and is strictly decreasing. Clipping its root to $[0,B]$
gives $z^*=\min\{B,\max(a/R_*-1,0)\}$, which is
\cref{eq:budgetcrit}. If $R_*=0$, positive local marginal production
uses the whole budget. Ties between best external channels affect the
composition of the remainder, not the local allocation.

\subsection{Increasing returns and an exceptional global test}
\label{app:returns-proof}
\begin{proof}[Proof of \cref{prop:exponential-necessity}]
The necessary logarithmic test, which does not require concavity of
$\ell$, implies $p_j>S_j\ell(B)/\ell'(B)$ wherever $S_j>0$.
Therefore $\ell'(B(t))/\ell(B(t))>S_j(t)/p_j(t)\geq s/M$
at every sufficiently late time. Choose such a time $t_0$.
Continuity and $B(t)\to\infty$ imply that every $B>B(t_0)$ is attained
at some $t\geq t_0$, by the intermediate value theorem. Thus the
strict differential inequality holds throughout that budget interval.
Integrating it on any $[B_0,B]$ contained there gives
$\ln\ell(B)-\ln\ell(B_0)>(s/M)(B-B_0)$.
For convex $\ell$ with $\ell(0)=0$, the tangent inequality at $B$
is $0\geq\ell(B)-B\ell'(B)$, so $\Gamma\leq B$.
Direct differentiation gives $\Gamma=B/k$ for $\ell(z)=z^k$ and
$\Gamma=[1-\exp(-\psi B)]/\psi$ for the exponential specification.
Every fixed finite $k$ still yields $\Gamma\to\infty$, contradicting
the strict price floor when prices are bounded and residual slopes stay
positive.
\end{proof}

\begin{lemma}[Exponential production: an exact endpoint test]
\label{lem:exponential-endpoints}
Let $B,a,p,\psi>0$, $d_0\geq1$, and $r,S\geq0$. On $[0,B/p]$,
define
\begin{equation}
 U(x)=\ln\!\left[d_0+\frac a\psi
   \{\exp[\psi(B-px)]-1\}+rx\right]+Sx.
 \label{eq:exponential-objective}
\end{equation}
Then $x=0$ is a global maximizer if and only if
$U'(0)\leq0$ and $U(0)\geq U(B/p)$.
The result allows ties and does not extend to arbitrary concave
residual utility or arbitrary concave external computation.
\end{lemma}
\begin{proof}
Necessity follows from a feasible right derivative and endpoint
comparison. For sufficiency, put $k=\psi p$,
$A=(a/\psi)\exp(\psi B)$, and $c=d_0-a/\psi$. The positive
argument of the logarithm is $D(x)=A\exp(-kx)+rx+c$.
The derivative has the sign of
\[
 f(x)=D(x)U'(x)=(S-k)A\exp(-kx)+Sr x+r+Sc.
\]
If $S\leq k$, then
$f'(x)=k(k-S)A\exp(-kx)+Sr\geq0$. Hence $f(0)\leq0$ makes
the non-positive set an initial interval. If $S>k$, then
$f''(x)=k^2(S-k)A\exp(-kx)>0$. Its non-positive sublevel set
is convex and contains zero, so is again an initial interval, possibly
a singleton or the entire domain. In both cases $U$ is non-increasing
and then non-decreasing. Its maximum is consequently attained at an
endpoint. The assumed endpoint comparison proves the result.
For $r=0$, $a=\psi d_0$, and $S=\psi p$, the function is constant;
thus a strict uniqueness conclusion would be false.
\end{proof}

\subsection{Smooth counterexamples to the endpoint shortcut}
\label{app:exponential-counterexamples}
These constructions use positive logarithm arguments throughout and
infinitely differentiable, non-decreasing, concave alternatives.
They refute extensions of \cref{lem:exponential-endpoints}, not its
stated affine and linear assumptions. An interior comparison suffices;
locating the precise global optimizer is unnecessary.

First set $B=10$, $p=\psi=d_0=1$, $a=0.01$, and $r=0$, and replace
$Sx$ by
\[
 V_{\rm ext}(x)=\frac{0.76}{20}
 \left\{\ln(1+\exp(140))-
             \ln(1+\exp[20(7-x)])\right\}.
\]
Here $V_{\rm ext}(0)=0$,
$V_{\rm ext}'(x)=0.76/[1+\exp(20(x-7))]>0$, and
$V_{\rm ext}''(x)<0$. Let
$U(x)=\ln[1+0.01\{\exp(10-x)-1\}]+V_{\rm ext}(x)$.
The following strict intervals are sufficient:
\begin{align*}
 -0.236&<U'(0)<-0.235, &5.399&<U(0)<5.400,\\
 5.468&<U(7)<5.469, &5.319&<U(10)<5.321.
\end{align*}
Thus the first-order condition and comparison with the all-external
endpoint both pass, but the interior value at $x=7$ is larger.
These bounds can be checked directly from the displayed expressions,
using $\ln(1+\exp v)=v+\ln(1+\exp(-v))$ for $v\geq0$
and $0<\ln(1+u)<u$ for $u>0$; the gaps exceed $0.06$.
A looser analytic certificate already suffices: $U(0)<5.4$,
$U(7)>0.17+0.038(140-0.694)>5.463$, and $U(10)<5.32<U(0)$.
These follow from $\exp(3)>20$, $\ln(1.19)>0.17$,
$\ln2<0.694$, $\exp(10)<22027$, $\exp(5.4)>221.4$, and
$\ln100<4.606$. Furthermore, the negative computational derivative
has magnitude greater than $0.99$ while the residual derivative is
less than $0.76$, so $U'(0)<0$. The numerical optimizer is an
additional check, not the argument for the counterexample.

Second set $B=3$, $p=\psi=a=d_0=1$, and $S=0$, and replace affine
external production by
\[
 G(x)=\frac9{50}\left\{\ln(1+\exp(100))-
                  \ln(1+\exp[50(2-x)])\right\}.
\]
Then $G(0)=0$, $G'(x)=9/[1+\exp(50(x-2))]>0$, and $G''(x)<0$.
For $U(x)=\ln[\exp(3-x)+G(x)]$, one has
\[
 U'(0)=-1+\frac{9}{\exp(3)[1+\exp(-100)]}<0,
 \qquad U(0)=3.
\]
Also $G(3)<18$, so $U(3)<\ln19<3$, whereas
$G(2)>18-(9/50)\ln2$ gives
\[
 \exp(1)+G(2)>\exp(1)+18-(9/50)\ln2>\exp(3),
\]
so $U(2)>3$. For example, the elementary bounds
$\exp(1)>2.718$, $\exp(3)<20.086$, and $\ln2<0.694$
already prove the last strict comparison. The two-test shortcut fails
even with zero residual utility when the external production function
is allowed to be concave rather than affine.

\subsection{Existence and uniqueness of the tipping budget}
\label{app:tipping-proof}
\begin{proof}[Proof of \cref{prop:tipping}]
The assumed positive derivative makes $\Gamma$ well defined and
continuously differentiable on $[0,\infty)$. The quotient rule gives
\cref{eq:gamma-monotone}. Hence $\Gamma$ starts at zero, is strictly
increasing, and is unbounded because $\Gamma(B)\geq B$.
If a positive slope exists, the finite number
$\min_{S_j>0}p_j/S_j$ is strictly positive. The intermediate value
theorem and strict monotonicity give exactly one positive crossing.

Write $d_j(B)=p_j\ell'(B)-S_j\ell(B)
=\ell'(B)[p_j-S_j\Gamma(B)]$.
Below the crossing each $d_j$ is positive; above or at it, a direction
attaining the minimum has $d_j\leq0$ and $r_j+d_0S_j>0$.
The exact logarithmic threshold then proves attainability below and
non-attainability at or above the crossing. On the lower interval,
$d_j'(B)=p_j\ell''(B)-S_j\ell'(B)\leq0$. Thus every nontrivial
ratio $(r_j+d_0S_j)/d_j(B)$, and their finite maximum, is
non-decreasing. At least one positive numerator has a denominator
tending to zero as $B\uparrow B_{\rm tip}$, proving divergence.

If $B(t)\leq B'<B_{\rm tip}$ eventually, the threshold is bounded
above by $a_{\rm crit}^{\ell}(B')<\infty$. The revised hypothesis
states that $a(t)$ strictly exceeds this finite bound eventually;
divergence of productivity is unnecessary. Every nontrivial corner inequality
is then strict; in a trivial direction it is also strict because
$p_j>0$ and the local marginal value is positive. The hub proposition
therefore makes every optimizer all-inward. If all $S_j=0$, all
thresholds are $r_j/[p_j\ell'(B)]$, finite at every finite budget.
There is then no residual tipping budget; the definition
$B_{\rm tip}=\infty$ records this case.
\end{proof}

\subsection{Value curvature at fixed and bounded budgets}
\label{app:curvature-proof}
\begin{proof}[Proof of \cref{prop:curvature}]
Fix $B>0$ and abbreviate $L_\ell=\ell(B)>0$ and
$d_\ell=\ell'(B)>0$. A direction passes exactly when
\[
 P_{j,\eta}(a):=\frac{r_j+S_j(d_0+aL_\ell)^\eta}{a d_\ell}
 \leq p_j.
\]
For a nontrivial direction and $0\leq\eta<1$,
\[
 P_{j,\eta}'(a)=\frac{-r_j+S_j(d_0+aL_\ell)^{\eta-1}
                    [(\eta-1)aL_\ell-d_0]}{a^2d_\ell}<0.
\]
The expression tends to infinity as $a\downarrow0$ and to zero as
$a\to\infty$. Continuity and strict monotonicity prove the unique
finite threshold. At $\eta=1$ expansion gives
\cref{eq:general-production} and hence \cref{eq:acrit-general}.
At $\eta>1$ and $S_j>0$, the sign condition is the non-positive
sublevel set of
\[
 r_j+S_j(d_0+aL_\ell)^\eta-p_ja d_\ell.
\]
This function is strictly convex, positive at zero, and tends to
infinity as $a\to\infty$. Its non-positive set is therefore empty,
a singleton, or a bounded closed interval separated from zero.
For $S_j=0$ the exact inequality is simply
$a\geq r_j/(p_jd_\ell)$, with a trivial direction unrestricted.
Putting $L_\ell=\ln(1+B)$ and $d_\ell=(1+B)^{-1}$ recovers
\cref{eq:iso-corner} and every fixed-budget statement of the original
curvature proposition.

For the bounded-budget conclusions, take a compact positive-budget
interval containing all sufficiently late $B(t)$ and the limit
$\bar B>0$. On this interval $L_\ell$ is bounded above and below by
positive constants and $d_\ell$ is likewise bounded away from zero
and above. These statements follow from the maintained differentiability
and concavity assumptions locally, or directly from the $C^2$ assumptions
of the tipping result. For $\eta<1$, the preceding required prices
are uniformly $O(a^{-1}+a^{\eta-1})$ and hence tend to zero as a
function of $a$. Choose a finite threshold above which these bounds are
strictly below every positive $p_j$; it suffices for the actual path to
eventually exceed that threshold. For $\eta>1$ and $S_j>0$, they are
uniformly bounded below by a positive constant times $a^{\eta-1}$.
A finite threshold above which this lower bound exceeds $p_j$ likewise
suffices for exclusion. At $\eta=1$, a limiting budget below the tipping
point is eventually bounded by some $B'<B_{\rm tip}$, and eventual
$a(t)>a_{\rm crit}^{\ell}(B')$ suffices. A limit above the tipping point
eventually lies above it.
At equality, the exact pointwise comparison with the diverging
threshold is indispensable.
\end{proof}

Two boundaries clarify the quantifiers. Mere upper boundedness of
$B(t)$ is insufficient for the high-curvature exclusion: with
$\ell(z)=z$, $\eta=2$, $r=0$, $S=p=d_0=1$, $B=a^{-2}$, the
corner condition is $(1+1/a)^2\leq a$, true for $a\geq4$.
At logarithmic value, take $\ell(z)=z$, $r=0$, $S=p=d_0=1$,
so $B_{\rm tip}=1$. Along $B(t)=1-\exp(-t)$, the threshold is
$a_{\rm crit}=\exp(t)$. The paths $a(t)=\exp(2t)$ and
$a(t)=\exp(t/2)$ give respectively inward and outward allocations
for $t>0$ despite having the same limiting budget.

\subsection{A complete inward-basin characterization}
\label{app:basin-proof}
\begin{proof}[Proof of \cref{prop:basin}]
For $z=B-px$, put $D=d_0+a\ln(1+z)+rx$. This is positive,
$D''(x)=-ap^2/(1+z)^2<0$, and therefore
$(\ln D+Sx)''=D''/D-(D'/D)^2<0$.
The static optimizer is unique. It depends continuously on $(B,a)$:
rescale $x=(B/p)v$ to the compact fixed interval $v\in[0,1]$;
continuity of the objective and uniqueness imply that every convergent
subsequence of optimizing $v$ converges to the unique limiting
optimizer.

The derivative has the sign of
\begin{equation}
 T(x;B,a)=r+Sd_0+Sr x+
       a\left[S\ln(1+z)-\frac{p}{1+z}\right].
 \label{eq:basin-derivative}
\end{equation}
At fixed feasible $x$, its partial derivative in $B$ is
$a[S/(1+z)+p/(1+z)^2]>0$. Increasing $B$ also enlarges the
feasible interval, so strict concavity implies that $x^*(B,a)$ is
non-decreasing in $B$. At an interior optimum,
\[
 S\ln(1+z)-p/(1+z)=-(r+Sd_0+Sr x)/a<0.
\]
Thus increasing $a$ makes the derivative at that old optimum strictly
negative, and the new optimizer is strictly smaller. At $x=0$, a
non-positive derivative implies the same negative bracket, so a larger
$a$ preserves the zero optimum. At the other endpoint a new optimizer
cannot exceed the feasible upper bound. These observations prove that
$x^*$ is non-increasing in $a$, and strictly decreasing whenever the
old optimizer is interior and the new optimizer remains positive.

For $B<B_{\rm tip}$, the zero derivative test is
$a\geq a_{\rm crit}^{\ell}(B)$. If this is met at some date, the
budget stays constant, productivity rises strictly, and all subsequent
choices remain zero. At or above the tipping budget, let
$x_\dagger=(B-B_{\rm tip})/p$. It is feasible and leaves local
investment $z=B_{\rm tip}$. The bracket in \cref{eq:basin-derivative}
vanishes there, while $r+Sd_0+Sr x_\dagger>0$.
Consequently $x^*>x_\dagger$, proving \cref{eq:runaway-bound}.
Reaching equality with the tipping budget is followed by a strict
crossing at the next date, after which the positive excess grows at
least geometrically with multiplier $1+g_X/p>1$.

Call a path captured if it chooses zero in finite time, and crossed if
it reaches or exceeds $B_{\rm tip}$ in finite time. A path that does
neither has $x_t>0$ and a non-decreasing budget bounded by
$B_{\rm tip}$. Hence $B_t\to B_\infty\leq B_{\rm tip}$.
If $B_\infty<B_{\rm tip}$, the static productivity threshold is
bounded along this path, whereas $a_t\to\infty$, forcing capture.
Thus every unclassified path must converge to $B_{\rm tip}$ from below.

The captured set of initial productivities is nonempty: it contains
$a_0\geq a_{\rm crit}^{\ell}(\bar B)$. The crossed set is also
nonempty. Choose a finite integer $N$ such that
$\bar B(1+g_X/p)^N>B_{\rm tip}$. Choose $a_0>0$ small enough that
$a_0\exp(g_at)\leq Sd_0/p$ for $0\leq t<N$. At the all-external
endpoint the derivative numerator is
$r+Sd_0+SrB/p-ap\geq0$, so the optimizer is all-external throughout
those periods. The resulting budget crosses by period $N$.

Both sets are open as subsets of $(0,\infty)$. A captured path has
$B<B_{\rm tip}$, and one period after capture its corner inequality
is strict because productivity has risen at unchanged budget. A
crossed path is strictly above the tipping budget no later than one
period after first reaching it. Continuity of every finite sequence of
states and decisions preserves these strict finite-time certificates
under sufficiently small changes of $a_0$.

The two sets are ordered. A larger $a_0$ gives smaller investment at a
common budget; induction gives weakly smaller budgets at every date.
If a smaller-$a_0$ path is captured, the larger one is captured no later
than that date, since $a_{\rm crit}^{\ell}$ is non-decreasing.
Similarly, crossing of a larger-$a_0$ path implies crossing of a smaller
one. Thus the captured set is an upper interval and the crossed set a
lower interval.

There is at most one initial productivity in neither set. Suppose
$a_0^-<a_0^+$ generated two such paths. Then $B_t^-\geq B_t^+$,
and the gap is non-decreasing, because
$x^*(B_t^-,a_t^-)\geq x^*(B_t^+,a_t^+)$. Both budgets would have to
converge to $B_{\rm tip}$, so their gap would have to remain identically
zero. Yet for all sufficiently late dates, neither optimum is the
all-external endpoint: at that endpoint its derivative numerator is
$r+Sd_0+SrB/p-ap$, which becomes negative for bounded $B$ and
$a\to\infty$. Neither optimum is zero by assumption. They are then
interior, and the strict productivity comparison at a common budget
forces a positive investment gap, a contradiction.

Finally, two disjoint nonempty open intervals cannot exhaust the
connected set $(0,\infty)$. Their boundary contains a point in neither
set, and the preceding argument makes that point unique. Denote it
$a_{\rm sep}$. The low-productivity construction makes it positive;
openness of the captured set around $a_{\rm crit}^{\ell}(\bar B)$
makes it strictly smaller than that threshold. The boundary path has
$B_t\uparrow B_{\rm tip}$, and telescoping
$B_{t+1}-B_t=g_Xx_t$ gives its stated total investment.

For comparative statics, an increase of $g_a$ increases every future
productivity and weakly decreases investment and budgets by induction.
An increase of $g_X$ weakly increases the transition
$B\mapsto B+g_Xx^*(B,a)$, and an increase of $\bar B$ increases its
initial state. Both therefore weakly increase all budgets. The ordered
capture and crossing sets imply the asserted monotonicities of their
unique separator. The comparison in $\bar B/B_{\rm tip}$ holds with
$p,S,r,d_0$ fixed, so that $B_{\rm tip}$ itself is fixed.
\end{proof}

\begin{lemma}[No-further-investment certificate on a mobilization path]
\label{lem:stay-certificate}
In the same one-channel logarithmic model, let a specified
$B_{\rm loc}(t)>0$ be non-decreasing, bounded by $\bar B$, and let
$a_{T+k}=a_T\exp(g_ak)$ with $g_a>0$. From state $(T,X_T)$,
zero investment at every date $T+k$ is optimal if and only if
\begin{equation}
 a_T\geq a_{\rm stay}(T,X_T):=
 \sup_{k\geq0}\exp(-g_ak)
 a_{\rm crit}^{\ell}\bigl(B_{\rm loc}(T+k)+g_XX_T\bigr),
 \label{eq:stay-certificate}
\end{equation}
where the threshold is extended to infinity at and above
$B_{\rm tip}$. In particular,
$\bar B+g_XX_T<B_{\rm tip}$ and
$a_T\geq a_{\rm crit}^{\ell}(\bar B+g_XX_T)$ are sufficient.
For the constant-$\bar B$ model starting at $X_0=0$, another sufficient
capture certificate is the existence of an integer $T\geq0$ such that
\begin{equation}
 \bar B(1+g_X/p)^T<B_{\rm tip},\qquad
 a_0\exp(g_aT)\geq
 a_{\rm crit}^{\ell}\bigl(\bar B(1+g_X/p)^T\bigr).
 \label{eq:finite-capture-certificate}
\end{equation}
\end{lemma}
\begin{proof}
If all future investments are zero, the stock is fixed at $X_T$ and
the static corner test must hold at every date. These tests are exactly
the inequalities whose supremum appears in \cref{eq:stay-certificate}.
Conversely, when they hold, induction gives the unique zero optimizer
and unchanged stock at every date. The upper-budget certificate uses
monotonicity of $a_{\rm crit}^{\ell}$ and rising productivity.
In the constant-ceiling model, $x_t\leq B_t/p$ implies
$B_t\leq\bar B(1+g_X/p)^t$. Under
\cref{eq:finite-capture-certificate}, the actual budget at date $T$
is below the tipping budget and its threshold does not exceed the
stated bound. Hence $x_T=0$, after which capture is permanent.
\end{proof}

The cases excluded from the strict basin theorem have direct answers.
If $g_X=0$ and $\bar B<B_{\rm tip}$, the fixed budget and $g_a>0$
ensure eventual capture for every $a_0>0$, but not necessarily immediate
inwardness. If $g_a=0$, $g_X>0$, and the initial budget is below the
tipping point, $a_0\geq a_{\rm crit}^{\ell}(\bar B)$ is captured
immediately. Below that threshold, investment is initially positive and
is non-decreasing as the budget grows; thus the tipping budget is crossed
in finite time. If $\bar B\geq B_{\rm tip}$, every finite productivity
is outward immediately. If $S=0$, the finite residual tipping budget and
the separator construction do not apply. At $B=0$, the feasible set is
a singleton and the divided corner formulas are unnecessary.

For the uniform bounded-budget extension at $\eta<1$, concavity gives
$\ell'(B)\geq\ell'(B_{\max})>0$ and $\ell(B)\leq\ell(B_{\max})$.
The required price is therefore at most
$[r_j+S_j(d_0+a\ell(B_{\max}))^\eta]/[a\ell'(B_{\max})]$,
which tends to zero. This proves the assertion even when budgets approach
zero. It does not prove the analogous $\eta>1$ statement without a positive
lower budget bound.


\subsection{Effective budgets with a finite productivity ceiling}
\label{app:ceiling-budget-proof}
\begin{proof}[Proof of \cref{cor:ceiling-budget}]
Let $f(B)=a_{\rm crit}^{\ell}(B)$ for $0<B<B_{\rm tip}$.
Every denominator in \cref{eq:acrit-general} is positive there.
The finite maximum of the resulting continuous ratios is continuous.
The proof of \cref{prop:tipping} establishes that $f$ is non-decreasing
and $f(B)\to\infty$ at $B_{\rm tip}$. Thus its sublevel set
$E=\{B\in(0,B_{\rm tip}):f(B)\leq a_{\max}\}$ is downward closed.
Divergence gives a neighbourhood below $B_{\rm tip}$ disjoint from $E$.
If $E$ is empty, the stated convention gives $B^*=0$, and no positive
budget is attainable under the ceiling. Otherwise its positive supremum
lies strictly below $B_{\rm tip}$. A sequence in $E$ tending to that
supremum, continuity, and downward closure prove $E=(0,B^*]$.
If $f(B^*)<a_{\max}$, continuity would put a larger budget in $E$;
hence $f(B^*)=a_{\max}$. The exact corner criterion now proves all
attainability and endpoint claims. A larger ceiling enlarges $E$.
For every fixed $b<B_{\rm tip}$, all ceilings at least $f(b)$ have
$B^*\geq b$. Together with $B^*<B_{\rm tip}$ this proves convergence.
The upper bound on $B_{\rm tip}$ follows from $\Gamma(B)\geq B$.

No strict monotonicity of $f$ was used. For example, take $\ell(B)=B$,
$d_0=1$, two unit prices, and $(r_1,S_1)=(0,1)$,
$(r_2,S_2)=(2,0)$. Then $B_{\rm tip}=1$ and
$f(B)=\max\{(1-B)^{-1},2\}$. A ceiling below two admits no positive
budget; at two, $B^*=1/2$ even though $f(0)=2$. This plateau also
warns against replacing existence by uniqueness of an inward optimizer.
At $a=2$, the second channel can exchange local and remote computation
without changing total computation or residual value. Strict inequalities,
not weak ceiling feasibility, justify the assertion about every optimizer
in the revised path statement of \cref{prop:tipping}.
\end{proof}

\subsection{Bounded productivity below unit value curvature}
\label{app:ceiling-growth-proof}
For any fixed $b_0>0$, concavity gives
$\ell(B)\leq\ell(b_0)+\ell'(b_0)(B-b_0)$ for $B\geq b_0$.
The finite derivative therefore supplies a finite $K>0$ with
$\ell(B)\leq KB$ eventually; differentiability at zero is unnecessary.
A ceiling gives $\Omega\leq a_{\max}KB$. For $0<\eta<1$, dividing
\cref{eq:necessary-output-growth} by $B$ requires a quantity bounded
above by $a_{\max}K$ to dominate a positive multiple of
$B^{\eta/(1-\eta)}$, a contradiction. More generally the same calculation
bounds any eventually inward productivity below by
$K^{-1}(s/M)^{1/(1-\eta)}B^{\eta/(1-\eta)}$, proving $a(t)\to\infty$.
At $\eta=0$, \cref{eq:linear-value-growth} is exact. Replacing actual
$a(t)$ by its upper bound yields only a necessary test. If its right-hand
side is bounded below by a positive constant and $\ell'(B(t))\to0$,
that test fails. Linear production, whose derivative does not vanish,
can instead satisfy it with a fixed finite productivity.

\subsection{The capped basin and its necessary qualification}
\label{app:capped-basin-proof}
\begin{proof}[Proof of \cref{cor:capped-basin}]
The continuity and monotonicities of $x^*(B,a)$ proved in
\cref{app:basin-proof} do not depend on the productivity path.
In case (i), put $\delta_x=x^*(\bar B,a_{\max})>0$.
Since $B_t\geq\bar B$ and $a_t\leq a_{\max}$, those monotonicities give
$x_t\geq\delta_x$ until crossing and
$B_t\geq\bar B+t g_X\delta_x$. Crossing therefore occurs no later than
$\lceil(B_{\rm tip}-\bar B)/(g_X\delta_x)\rceil$.
The runaway bound is a static inequality for every finite $a>0$, so it
continues to apply under either ceiling path.

For case (ii), the one-channel logarithmic threshold
\[
 f(B)=\frac{r+Sd_0}{p/(1+B)-S\ln(1+B)}
\]
is strictly increasing below $B_{\rm tip}$. Write $b_*=B^*(a_{\max})$.
Then $\bar B<b_*<B_{\rm tip}$ and $f(b_*)=a_{\max}$.
As before, call a path captured when some $x_t=0$, and crossed when it
reaches $B_{\rm tip}$. Capture is permanent because productivity is
non-decreasing and the budget then stops changing. Every captured path
ends below $b_*$: a path with $a_0<a_{\max}$ has $a_t<a_{\max}$ at
all finite dates, while $a_0=a_{\max}$ captures at $\bar B<b_*$.

An uncaptured, uncrossed path has positive investments and bounded
increasing budgets, so $B_t\to B_\infty\leq B_{\rm tip}$ and $x_t\to0$.
Continuity and $a_t\to a_{\max}$ imply
$x^*(B_\infty,a_{\max})=0$. The positive derivative at $B_{\rm tip}$
rules out that endpoint, and the corner test gives $B_\infty\leq b_*$.
At every finite date noncapture gives $a_t<f(B_t)$. Continuity below
the tipping budget then gives $a_{\max}\leq f(B_\infty)$, so
$B_\infty\geq b_*$. Hence $B_\infty=b_*$, and telescoping yields
$g_X\sum_{t\geq0}x_t=b_*-\bar B$.

The captured initial values form an upper interval and crossed values a
lower interval, by the policy monotonicities and pathwise ordering.
The captured interval is nonempty since it contains $a_0\geq f(\bar B)$;
the crossed interval is nonempty by the additional hypothesis. Both are
relatively open in $(0,a_{\max})$: one period after capture the corner
inequality is strict because $a_0<a_{\max}$, while a crossing is strict
at most one period after reaching $B_{\rm tip}$. Finite-horizon
continuity preserves these certificates. Connectedness supplies an initial
value in neither interval.

There cannot be two such values. If $a_0^-<a_0^+$, their productivities
remain strictly ordered at every finite date under
\cref{eq:saturating-productivity}. The corresponding budget gap is
non-negative and non-decreasing. If both limits were $b_*$, this gap
would be identically zero. Since $x_t\to0$ and $B_t\geq\bar B>0$,
their positive optima are eventually interior; strict productivity
monotonicity at the common budget then makes their investments different,
a contradiction. The unique remaining value is $a_{\rm sep}$.
Openness of capture at $f(\bar B)<a_{\max}$ and the existence of a
positive crossed value give $0<a_{\rm sep}<f(\bar B)$.
This proves the classification with no claim that every strictly ordered
family supplies a nonempty crossed basin.
\end{proof}

\paragraph{Counterexample without a crossed path.}
Let $p=S=d_0=g_X=1$, $r=0$, $\bar B=1/10$, $a_{\max}=10$,
and $g_a=\ln2$. For every $0<a_0\leq10$, the first budget satisfies
$B_1\leq1/5$, while $a_1=5+a_0/2>5$.
Using $\ln(6/5)<1/5$ gives
\[
 f(1/5)=\frac1{5/6-\ln(6/5)}<\frac{30}{19}<5.
\]
Consequently all paths are captured by date one, although
$a_{\max}>f(\bar B)$ and the family preserves strict order at every
finite date. Thus strict order alone does not establish a positive
separator with both basins nonempty.

\paragraph{A hard cap changes the boundary classification.}
Keep $p=S=d_0=g_X=1$, $r=0$, and $\bar B=1/10$, but set
$a_{\max}=f(1/5)$ and $a_t=\min(10^ta_0,a_{\max})$.
For every $a_0\in[a_{\max}/10,1]$, the date-zero all-external endpoint
is optimal because its derivative numerator is $1-a_0\geq0$.
Thus $B_1=1/5$; at date one productivity reaches the cap and the corner
condition is an equality. Strict concavity makes $x_1=0$, and the path
stays at $B_t=B^*(a_{\max})=1/5$ thereafter.
For $0<a_0<a_{\max}/10$, the first budget is again $1/5$, but
$x_1>0$ and $B_2>B^*$. Once above $B^*$, the uniform lower bound
$x^*(B_2,a_{\max})>0$ forces finite crossing of $B_{\rm tip}$.
The dividing initial value is still unique in this example, but its
path is captured in finite time, and a whole interval of initial values
shares the same terminal boundary budget. The unique \emph{uncaptured}
critical-path conclusion cannot be transferred to this coalescing family.

\subsection{Shared technological change}
\label{app:shared-technology-proof}
\begin{proof}[Proof of \cref{cor:shared-technology}]
The fixed $B,d_0,S_j$ leave $H_j=(1+B)\ln(1+B)S_j$ unchanged.
The transformed $A_j$ is $(1+B)(\lambda_r r_j+d_0S_j)$.
Divide these by $p_j/\kappa$ and by
$(p_j/\kappa)\lambda_a a$, respectively, and subtract the original
remote-return term to obtain \cref{eq:shared-technology}.
No division by $r_j$, $S_j$, or $A_j$ is involved, so zero slopes cause
no exception. If $\lambda_r=\lambda_a$, the remote-computation component
of $\zeta_j$ changes only by $\kappa$, whereas its residual component
changes by $\kappa/\lambda_a$. In either case $\xi_j$ changes by
$\kappa$, independently of productivity. In particular,
$\kappa\xi_j\geq1$ in a nontrivial residual direction rules out inwardness
at every finite transformed productivity. This does not assert that a
particular physical technology realizes these multipliers.
\end{proof}
